% PAGES: 235-236
% PDF page 235 (folio 219) opens with the "7.7 Multivariate random variables
% formulation for covariance and correlation matrices" section heading at
% the very top of the page; PDF page 234 (folio 218, the last page of chunk
% c3's range) was rendered and inspected directly and ends cleanly with
% Lemma 7.9's proof closed by its QED box, with blank space below -- so
% nothing is open to close at the start of this file.

\section{Multivariate random variables formulation for covariance and correlation matrices}

The formulas from section 7.1 connecting the covariance and correlation matrix of
random variables, as well as the Linear Transformation Property from section 7.3, can
be established in an elegant way by using multivariate random variables.

Let $X_1, X_2, \ldots, X_n$ be random variables on the same probability space, and let
$\bfX$ be the $n$-dimensional multivariate random variable
\begin{equation}
\bfX \;=\; \begin{pmatrix} X_1 \\ X_2 \\ \vdots \\ X_n \end{pmatrix}.
\end{equation}

The expected value $\mu_{\bfx}$ of $\bfX$ is
\[
\mu_{\bfx} \;=\; \E{\bfX} \;=\; \begin{pmatrix} \E{X_1} \\ \E{X_2} \\ \vdots \\ \E{X_n}
\end{pmatrix} \;=\; \begin{pmatrix} \mu_1 \\ \mu_2 \\ \vdots \\ \mu_n \end{pmatrix},
\]
where $\mu_i = \E{X_i}$ for all $i = 1:n$.

Similarly, if $\bfX = (X_{j,k})_{1\le j\le m,1\le k\le n}$ is a multivariate random
variable taking values in $\R^{m\times n}$, where $X_{j,k}$ are univariate random
variables on the same probability space, for all $1\le j\le m$, $1\le k\le n$, then
$\E{\bfX} = (\E{X_{j,k}})_{1\le j\le m,1\le k\le n}$.

The general result from Lemma 7.10 follows from the linearity of expected values of
random variables and will be used repeatedly:

\begin{lemma}
Let $\bfX$ be an $m \times n$ multivariate random variable. Let $A$ and $B$ be matrices
of sizes $p \times m$ and $n \times q$, respectively, with constant real entries. Then,
\begin{equation}
\E{A\bfX B} \;=\; A\,\E{\bfX}\,B.
\end{equation}
\end{lemma}

The result below is comprised of particular cases of (7.127) which will be used
subsequently.

\begin{lemma}
Let $\bfX$ be an $n$-dimensional multivariate random variable. Let $M$ be an $m\times n$
matrix with constant entries and let $C$ be an $n\times 1$ column vector with constant
entries. Then,\footnote{In formula (7.130), $\bfX C^t$ is an $n \times n$ matrix since
$\bfX$ is an $n \times 1$ column vector and $C^t$ is an $1 \times n$ row vector; cf.
(1.6).}
\begin{align}
\E{M\bfX} &= M\,\E{\bfX}; \\
\E{C^t\bfX} &= C^t\E{\bfX}; \\
\E{\bfX C^t} &= \E{\bfX}\,C^t.
\end{align}
\end{lemma}

The covariance matrix of the random variables $X_1, X_2, \ldots, X_n$ can be expressed
compactly using the multivariate random variable $\bfX$ given by (7.126) as follows:

\begin{lemma}
Let $\Sigma_{\bfx}$ be the covariance matrix of the random variables $X_1, X_2, \ldots,
X_n$. Then,
\begin{equation}
\Sigma_{\bfx} \;=\; \E{(\bfX - \E{\bfX})(\bfX - \E{\bfX})^t},
\end{equation}
where $\bfX = (X_i)_{i=1:n}$ is the corresponding multivariate random variable.
\end{lemma}

\begin{proof}
Note that $\bfX - \E{\bfX}$ is an $n \times 1$ column vector and $(\bfX - \E{\bfX})^t$ is
an $1 \times n$ row vector. Then, from (1.6), it follows that $(\bfX - \E{\bfX})(\bfX -
\E{\bfX})^t$ is the $n \times n$ matrix whose $(j,k)$ entry is $(X_j - \mu_j)(X_k -
\mu_k)$, for all $1 \le j,k \le n$. Thus, the $(j,k)$-th entry of the matrix
$\E{(\bfX-\E{\bfX})(\bfX-\E{\bfX})^t}$ is
\begin{equation}
\E{(X_j-\mu_j)(X_k-\mu_k)} \;=\; \cov(X_j,X_k),
\end{equation}
which is the same as the $(j,k)$-th entry of the covariance matrix $\Sigma_{\bfx}$; see
Definition 7.1.

We conclude that $\Sigma_{\bfx} = \E{(\bfX-\E{\bfX})(\bfX-\E{\bfX})^t}$.
\end{proof}

The formulas for the variance and the covariance of linear combinations of random
% Printed as "multivariate random variable" (singular) on PDF page 236 (folio
% 220); transcribed as printed, cf. fidelity rule.
variables from Lemma 7.3 can be proved by using multivariate random variable and the
formula (7.131) for the covariance matrix as seen below:

\begin{lemma}[Same as Lemma 7.3.]
Let $X_1, X_2, \ldots, X_n$ be random variables on the same probability space. Then,
\begin{equation}
\cov\!\left(\sum_{i=1}^n c_i^{(1)}X_i, \sum_{i=1}^n c_i^{(2)}X_i\right) \;=\;
\left(C^{(1)}\right)^t\Sigma_{\bfx}C^{(2)} \;=\; \left(C^{(2)}\right)^t\Sigma_{\bfx}C^{(1)},
\end{equation}
where $C^{(1)} = \left(c_i^{(1)}\right)_{i=1:n}$ and $C^{(2)} =
\left(c_i^{(2)}\right)_{i=1:n}$ are two column vectors of size $n$ with real entries,
and $\Sigma_{\bfx}$ is the covariance matrix of the random variables $X_1, X_2, \ldots,
X_n$.
\end{lemma}

\begin{proof}
Let $\bfX = (X_i)_{i=1:n}$ be an $n \times 1$ multivariate random variable. Note that
\[
\sum_{i=1}^n c_i^{(1)}X_i \;=\; \left(C^{(1)}\right)^t\bfX \quad \text{and} \quad
\sum_{i=1}^n c_i^{(2)}X_i \;=\; \left(C^{(2)}\right)^t\bfX.
\]

Then, using (7.132), we obtain that
\begin{align}
\cov\!\left(\sum_{i=1}^n c_i^{(1)}X_i, \sum_{i=1}^n c_i^{(2)}X_i\right)
&= \cov\!\left(\left(C^{(1)}\right)^t\bfX, \left(C^{(2)}\right)^t\bfX\right) \notag \\
&= E\!\left[\left(\left(C^{(1)}\right)^t\bfX - \E{\left(C^{(1)}\right)^t\bfX}\right)\right.\notag\\
&\qquad\qquad \left.\left(\left(C^{(2)}\right)^t\bfX - \E{\left(C^{(2)}\right)^t\bfX}\right)^t\right].
\end{align}

Using (7.129), we find that
\begin{align}
\left(C^{(1)}\right)^t\bfX - \E{\left(C^{(1)}\right)^t\bfX} &=
\left(C^{(1)}\right)^t\bfX - \left(C^{(1)}\right)^t\E{\bfX} \;=\;
\left(C^{(1)}\right)^t(\bfX - \E{\bfX}); \\
\left(\left(C^{(2)}\right)^t\bfX - \E{\left(C^{(2)}\right)^t\bfX}\right)^t &=
\left(\left(C^{(2)}\right)^t(\bfX-\E{\bfX})\right)^t \;=\; (\bfX-\E{\bfX})^tC^{(2)}.
\end{align}

From (7.134--7.136) and using (7.127), we conclude that
\begin{align*}
\cov\!\left(\sum_{i=1}^n c_i^{(1)}X_i, \sum_{i=1}^n c_i^{(2)}X_i\right)
&= \E{\left(C^{(1)}\right)^t(\bfX-\E{\bfX})(\bfX-\E{\bfX})^tC^{(2)}} \\
&= \left(C^{(1)}\right)^t \E{(\bfX-\E{\bfX})(\bfX-\E{\bfX})^t}\, C^{(2)} \\
&= \left(C^{(1)}\right)^t\Sigma_{\bfx}C^{(2)},
\end{align*}
% PDF page 236 (folio 220) ends here, immediately after this unnumbered
% align* block (three-line derivation with no printed equation tag). The
% closing "since Sigma_X = E[...]; see (7.131)." clause that finishes this
% sentence begins the top of PDF page 237 (folio 221) and continues in
% sec07_c4_2.tex -- the proof environment itself is NOT closed here, it
% remains open across this file boundary and is closed at the start of the
% next file.
